\documentclass[a4paper]{article}
\usepackage[pdftex]{graphicx}
\usepackage[utf8]{inputenc}
\usepackage{enumerate}
\usepackage{siunitx}
\sisetup{locale=DE}
\usepackage{href-ul}
\hypersetup{
	colorlinks=true,
	linkcolor=blue,
	urlcolor=blue}
\usepackage{icomma}
\usepackage{amssymb}
\hyphenation{Ver-simplification}
\usepackage{geometry}
\geometry{a4paper, top=15mm, left=15mm, right=15mm, bottom=15mm,
	headsep=10mm, footskip=12mm}

\begin{document}
	%Title
	{\bf Impromptu task from physics }\\
	\begin{enumerate}[1.]
		%Task 1
		{\bf \item What is the law of conservation of momentum?}
		%Exercise 2
		{\bf \item A minibus of mass $\SI{1.7}{\tonne}$ collides head-on with a small car of mass $ at speed $\SI{54}{\kilo\meter\over \hour}$ \SI{650}{\kilogram}$ together,} which was moving with the speed $\SI{36}{\kilo\meter\over \hour}$ before the collision. For simplicity, the collision should be assumed to be completely inelastic.
		
		\begin{enumerate}[a)]
			\item Calculate the speed of the two cars after the collision and indicate in which direction the two cars are moving. [wrong replacement result: $ u= \SI{27}{\kilo\meter\over \hour}$]
			\item What percentage of the kinetic energy existing before the accident is released by the impact
			Deformation work and internal energy converted?
		\end{enumerate}
		
		{\bf \item In the V2 rocket, the combustion gases flowed out at $\SI{2.0}{\kilo\meter\over \second}$.} The rocket thus developed a thrust of $\SI{2.0e5 }{\newton}$.
		
		\begin{enumerate}[a)]
			\item How much fuel had to be burned per second?
			\item The mass of the fully fueled rocket was $\SI{10}{\tonne}$. With what acceleration did the rocket lift off from the launch site in a vertical direction?
		\end{enumerate}
		
		{\bf \item Recoil: During the $\alpha$ decay of a uranium atom (m = 238 u), an $\alpha$ particle (m = 4 u) is ejected from the atomic nucleus.}
		It reaches a speed of $\SI{1.4e7}{\meter\over \second}$
		\begin{enumerate}[a)]
			\item What is the speed at which the remaining nucleus is repelled?
			\item What is the ratio of the two kinetic energies?
		\end{enumerate}
		
		
		{\bf \item Deep Impact: The asteroid that hit the Nördlinger Ries 15 million years ago had a mass of approx. $\SI{7.0e13}{\kilogram}$ and a speed of
			$\SI{20}{\kilo\meter\over \second}$ relative to the Earth.} Calculate the momentum it transferred to the Earth and decide whether the Earth could be thrown off its orbit by such an asteroid (mass Earth: $\SI{5.97e24}{\kilogram}$).
		
		\noindent
		\begin{minipage}{0.5\textwidth}
			{\bf \item Atwood falling machine}
			The masses $m_1=\SI{0.50}{\kilogram}$ and $m_2=\SI{0.40}{\kilogram}$ are connected via a cord and a pulley (both massless). The body $m_2$ is released from its position P1. The string has a loop and is only tensioned when the mass $m_2$ has fallen freely to a depth of $\SI{0.30}{\meter}$.
			\begin{enumerate}[a)]
				\item Calculate the speed at which the mass $m_1$ is then raised.
				\item What height $h_1$ does the mass $m_1$ reach?
			\end{enumerate}
		\end{minipage}
		\hfill
		\begin{minipage}{0.5\textwidth}
			\includegraphics[width=7 cm]{atwo136}
		\end{minipage}
		
		
		
	\end{enumerate}
	
	\fbox{
		\begin{minipage}{0.5\textwidth}
			For the solution, please \href{https://www.okuyakl.de/phys/p10impL136/le136.pdf}{click here} or scan the QR code.\\
			Additional worksheets are available at
			
			\href{http://www.okuyakl.com}{www.okuyakl.com}
		\end{minipage}
		\hfill
		\begin{minipage}{0.4\textwidth}
			\includegraphics[width=1.5 cm]{../../viecher/zwe03}
			\includegraphics[width=3 cm]{qre136}
			\includegraphics[width=2 cm]{../../viecher/afanticon1}
			
	\end{minipage}}
	
\end{document}